\chapter{解谜：向18种经典谜题的巅峰发起挑战}
\Large\textbf{The Puzzler: One Man's Quest to Solve the Most Baffling Puzzles, From Crosswords to Jigsaws to the Meaning of Life}
\par \emph{A. J. Jacobs}\normalsize

\par \textbf{纵横字谜}(crossword): 世界上已知最早的纵横字谜由Arthur Wynne设计, 1913年发表在The New York World. 《纽约时报》最初对纵横字谜持反对态度，1942年由编辑Margaret P. Farrar开始刊登纵横字谜；近三十年的纵横字谜编辑是Will Shortz\footnote{Will Shortz的大学毕业论文是关于美国早期的字谜历史. }. 字谜应当避免出现单词重复，或交叉的两个词都相对生僻的情况。藏词游戏(cryptics)主要流行于英国, 玩法与纵横字谜类似, 但线索与答案的联系更为复杂，一部分是同义，另一部分可能包括同音、回文等形式。1942年《每日电讯报》曾用藏词游戏选拔密码破译员进入布莱切利庄园。

\par \textbf{魔方}(Rubik's cube).

\par \textbf{回文字谜}(anagram): 将一个单词或短语字母重新排列得到另一个含义不同的单词或短语. 至少已有数百年历史, 伽利略使用过, 法王路易十三曾任命“王室回文创作师”. 回文是一种伟大的艺术(anagrams $\to$ ars magna). 始于1883年的美国解谜者联盟(National Puzzlers' League)在会刊The Enigma中有大量回文创作. 

\par 2016年Will Shortz创造了拼字蜜蜂(Spelling Bee)游戏, 规则是从7个六边形阵列的字母中用至少4个拼出单词, 且必须用到中心字母, 字母可以重复使用. 

\par \textbf{画谜}(rebus): 使用字母、插图、创意排版等手段隐藏一个单词或短语. 

\par \textbf{拼图}(jigsaw): 1760年前后英国制图师John Spilsbury将地图粘在木板上并按国家切成块, 作为地理教具, 这被认为是最早的拼图. 

\par \textbf{迷宫}(maze): 不同于迷园(labyrinth), 迷园没有分叉道路，从起点到终点无需做选择; 古希腊神话中囚禁牛头怪Minotaur的迷园其实是迷宫. 中世纪以来欧洲贵族有建造树篱迷宫的风气. Adrian Fisher是当今最高产的迷宫设计师之一, 作品包括迪拜高楼外立面迷宫等. 
\par 迷宫的解法: (1) Wall Follower: 右手放在墙上, 每到路口右转. 不适用于有墙与外墙不相连的情形. (2) Trémaux法: 标记每个路段的起点和终点, 避开有两个标记的路段. (3) Random Mouse: 每到路口随机决定方向. 

\par \textbf{数学与逻辑谜题}(math and logic puzzles): 可追溯到约公元前1500年古埃及的Rhind纸草书(Rhind mathematical papyrus), 包括87个涉及算术、代数、几何的数学问题. 神圣罗马帝国的查理曼大帝聘请Alcuin of York为其谜题设计师, 后者著有\textit{Problems to Sharpen the Young}. 1957年起, Martin Gardner在Scientific American上撰写专栏介绍数学谜题; 为纪念他, 数学谜题盛会G4G(Gathering for Gardner)每两年举办一次. 

\par \textbf{暗语和密码}(ciphers and secret codes): 暗语是词语或短语层面的替换，密码是字母层面的替换。频率分析是破译密码的最主要手段。历史上的密码包括15世纪的伏尼契手稿，黄道十二宫杀手的密信，Beatrix Potter的日记等。1990年James Sanborn为CIA设计的雕塑Kryptos上包含四段密码，至今仍有一段(K4)未被破译；已破译的三段使用了Vigenère密码和换位密码。

\par \textbf{视觉谜题}(visual puzzles): 包括找物型（找不同、找隐藏物体）、视错觉型等。历史上的谜题包括Pieter Bruegal的画作\emph{Netherlandish Proverbs}, Martin Handford的系列图书\emph{Where's Wally?}. Kit Williams于1979年出版的图画书\emph{Masquerade}包含了英格兰某处埋藏黄金兔子雕塑的线索，一度引发寻宝热潮。类似作品还有Byron Preiss的\emph{The Secret}, 用密文诗和画隐藏了美国、加拿大的12处宝藏。

\par \textbf{纸笔谜题}: 数独(Sudoku)最早由Howard Garns于1979年发明，当时称为Number Place. 1984年{\CJKfontspec{simsun.ttc}鍜}治真起将其改良并命名为数独。2004年宫本哲也发明了聪明格(KenKen). 

\par \textbf{国际象棋谜题}(chess puzzles): Genrikh Kasparyan的棋谜以困难著称. 西洋穴怪(Grotesques)指在子力明显处于劣势的情形下取胜的棋谜. 

\par \textbf{谜语}(riddles): 涉及谜语的作品包括《俄狄浦斯王》（斯芬克斯的谜语“什么东西早上四条腿，正午两条腿，晚上三条腿？”，谜底是人）；《爱丽丝漫游奇境》（疯帽子的谜语“乌鸦为什么像书桌？”）；《蝙蝠侠》（谜语人Riddler的谜语）；《霍比特人》；《哈利·波特与火焰杯》等。Charade指将谜底音节分开分别猜测的谜语。

\par \textbf{日本机关盒}(Japanese puzzle boxes): 龟井明夫开创了现代机关盒的时代，将机关盒的复杂程度和可玩性提升到新高度。美国机关盒设计师Kagen Sound曾设计一个包含22个机关谜题的书桌。

\par \textbf{争议谜题}: Monty Hall问题（“三扇门”）曾引起争议。睡美人问题(sleeping beauty problem)则至今仍无定论，主要观点为1/2, 1/3两派。

\par \textbf{解谜大赛}: MIT Puzzler Hunt创立于1981年，以团队形式解决系列谜题。
